% ============================================================================
% SYMBOLES DE LIAISONS — BIBLIOTHÈQUE CINÉMATIQUE
% Les commandes de cinematique.sty attendent les coordonnées sans parenthèses
% dans leur troisième argument : {0,0} et non {(0,0)}.
% ============================================================================

\renewcommand{\TLSymPivot}{%
  \begin{tikzpicture}[baseline=-.5ex,scale=.65]
    \scPivotx{L}{0,0}
  \end{tikzpicture}%
}

\renewcommand{\TLSymPivotGlissant}{%
  \begin{tikzpicture}[baseline=-.5ex,scale=.65]
    \scPivotGlissantx{L}{0,0}
  \end{tikzpicture}%
}

\renewcommand{\TLSymGlissiere}{%
  \begin{tikzpicture}[baseline=-.5ex,scale=.65]
    \scGlissierex{L}{0,0}
  \end{tikzpicture}%
}

\renewcommand{\TLSymHelicoidale}{%
  \begin{tikzpicture}[baseline=-.5ex,scale=.65]
    \scHelicoidalex{L}{0,0}
  \end{tikzpicture}%
}

\renewcommand{\TLSymAppuiPlan}{%
  \begin{tikzpicture}[baseline=-.5ex,scale=.70]
    \scAppuiPlan{L}{0,0}
  \end{tikzpicture}%
}

\renewcommand{\TLSymSpherePlan}{%
  \begin{tikzpicture}[baseline=-.5ex,scale=.70]
    \scSpherePlan{L}{0,0}
  \end{tikzpicture}%
}

\renewcommand{\TLSymCylindrePlan}{%
  \begin{tikzpicture}[baseline=-.5ex,scale=.70]
    \scCylindrePlanx{L}{0,0}
  \end{tikzpicture}%
}

\renewcommand{\TLSymSphereCylindre}{%
  \begin{tikzpicture}[baseline=-.5ex,scale=.70]
    \scSphereCylindrex{L}{0,0}
  \end{tikzpicture}%
}

\renewcommand{\TLSymSpherique}{%
  \begin{tikzpicture}[baseline=-.5ex,scale=.70]
    \scSpherique{L}{0,0}
  \end{tikzpicture}%
}

\renewcommand{\TLSymSpheriqueDoigt}{%
  \begin{tikzpicture}[baseline=-.5ex,scale=.70]
    \scSpheriqueDoigt{L}{0,0}
  \end{tikzpicture}%
}
